%{{{ preamble
\documentclass[a4paper,9pt]{article}
\usepackage{anysize}
\marginsize{2cm}{2cm}{1cm}{1cm}
%\textwidth 6.0in \textheight = 664pt
\usepackage{xltxtra}
\usepackage{xunicode}
\usepackage{graphicx}
\usepackage{color}
\usepackage{xgreek}
\usepackage{fancyvrb}
\usepackage{minted}
\usepackage{listings}
\usepackage{enumitem} 
\usepackage{framed} 
\usepackage{relsize}
\usepackage{float} 
\usepackage{pstricks}
\usepackage{pst-node}
\usepackage{pst-blur}
\setmainfont[Mapping=tex-text]{FreeSerif}
%}}}
\begin{document}

\def\thesection {\Roman{section}.}

\begin{titlepage}
\begin{center}
\begin{figure}[h] 
     \includegraphics[width=0.2\textwidth]{title/ntua_logo}
\end{figure}
\vspace{1cm}
\begin{LARGE}\textbf{ΕΘΝΙΚΟ ΜΕΤΣΟΒΙΟ ΠΟΛΥΤΕΧΝΕΙΟ\\[1.5cm]}\end{LARGE}
\begin{Large}
ΣΧΟΛΗ ΗΜ\&ΜΥ\\
Δίκτυα Επικοινωνιών\\[2cm]
4\textsuperscript{η} Άσκηση\\
Ακ. έτος 2011-2012\\
\end{Large}
\vfill
\begin{flushright}
\begin{tabular}{l r}
{Γρηγόρης \textsc{Λύρας}}&
{Α.Μ.: 03109687}\\
\end{tabular}
\end{flushright}

\large\today\\
\end{center}
\end{titlepage}




\section*{Aνάλθση αποτελεσμάτων με τη βοήθεια του NAM}
Αρχικά χρησιμοποιήσαμε τον κώδικα που μας δίνεται από την εκφώνηση:

\inputminted[fontsize=\footnotesize]{tcl}{files/ex1_1.tcl}


Τρέχουμε το animation και παρατηρούμε την εξέλιξη της πορείας μετάδοσης των
πακέτων:

\begin{itemize}
    \item Με τη βοήθεια του animation παρατηρούμε πως η αποστολή των 50
    πακέτων από τον κόμβο 0 στον κόμβο 3 (Go back N protocol window 7) ολοκληρώνεται τη
    στιγμή 1.3sec. Από την άλλη πλευρά το τελευταίο πακέτο από τον κόμβο 1
    στον κόμβο 2 (Stop and Wait protocol window 1) φτάνει τη χρονική στιγμη
    5.6sec.
    \item $\frac{2Mb}{1040*8bits}\aprox240packets$
\end{itemize}

\end{document}

